\documentclass[11pt,a4paper]{article}
\usepackage[T1]{fontenc}
\usepackage[utf8]{inputenc}
\usepackage{mathptmx}
\AtBeginDocument{\let\hbar\relax
  \DeclareMathSymbol{\hbar}{\mathord}{AMSb}{"7E}}
\usepackage{amsmath,amssymb,amsthm}
\usepackage{amsfonts}
\usepackage{geometry}
\usepackage{microtype}
\usepackage{hyperref}
\usepackage{booktabs}
\usepackage{enumitem}
\usepackage{natbib}
\usepackage{titlesec}

\geometry{top=1.1in,bottom=1.1in,left=1.15in,right=1.15in}

\titleformat{\section}{\normalfont\bfseries\large}{\thesection}{1em}{}
\titleformat{\subsection}{\normalfont\bfseries\normalsize}{\thesubsection}{1em}{}
\titleformat{\subsubsection}{\normalfont\bfseries\itshape\normalsize}{\thesubsubsection}{1em}{}
\titlespacing*{\section}{0pt}{14pt}{4pt}
\titlespacing*{\subsection}{0pt}{10pt}{3pt}
\titlespacing*{\subsubsection}{0pt}{8pt}{2pt}

\theoremstyle{plain}
\newtheorem{theorem}{Theorem}[section]
\newtheorem{proposition}[theorem]{Proposition}
\newtheorem{corollary}[theorem]{Corollary}
\theoremstyle{definition}
\newtheorem{definition}[theorem]{Definition}
\newtheorem{conjecture}[theorem]{Conjecture}
\theoremstyle{remark}
\newtheorem{remark}[theorem]{Remark}

\hypersetup{colorlinks=true,linkcolor=black,citecolor=black,urlcolor=black,
  pdftitle={Relational Actualism: Dark Matter as Actualization-Density Topology (RADM)},
  pdfauthor={Joshua F. Sandeman}}

\newcommand{\Pact}{\mathcal{P}_{\mathrm{act}}}
\newcommand{\Tact}{T^{\mu\nu}_{\mathrm{RA}}}
\newcommand{\ARA}{A_{\mathrm{RA}}}
\newcommand{\GG}{\mathcal{G}}
\newcommand{\AP}{\mathrm{AP}}
\newcommand{\MB}{\mathrm{MB}}
\newcommand{\ket}[1]{|#1\rangle}

% ─────────────────────────────────────────────────────────────────────────────
\begin{document}

\begin{titlepage}
\centering
\vspace*{1cm}
{\LARGE\bfseries Relational Actualism: Dark Matter as\\
Actualization-Density Topology ({RADM}):\par}
\vspace{0.4cm}
{\large\itshape Flat Rotation Curves, the Bullet Cluster, and the\\
Topological Boundary Mechanism\par}
\vspace{1.2cm}
{\large Joshua F.\ Sandeman\par}
\vspace{0.3cm}
{\normalsize\itshape Independent Researcher, Salem, Oregon\par}
\vspace{0.3cm}
{\normalsize March 2026\par}
\vspace{0.3cm}
  {\small\textit{Preprint:}~\href{https://doi.org/10.5281/zenodo.19198513}
    {doi.org/10.5281/zenodo.19198513}\par}
\vspace{1.5cm}

\begin{abstract}
The companion paper RAGC~\citep{Sandeman2026b} establishes that causally silent
dark matter candidates (WIMPs, axions, sterile neutrinos) cannot source macroscopic
spacetime curvature in Relational Actualism (RA): a 100~GeV WIMP generates
$10^{25}$--$10^{38}$ times less actualization flux than equivalent baryonic matter,
falling decisively below the threshold for gravitational lensing. This paper
develops the mechanisms by which the same actualization-density field accounts for
the observational signatures currently attributed to dark matter, without invoking
any new particle species. Derived as the static, weak-field limit of the covariant
actualization tensor established in RAGC~\citep{Sandeman2026b}, we utilise the
two-component metric sourcing equation $\nabla^2\Phi = 4\pi G[\rho_A + \rho_\lambda]$,
where $\rho_A$ tracks accumulated causal depth (recovering standard GR in the
equilibrium limit) and $\rho_\lambda \propto \nabla^2(\ln\lambda)$ is a gradient
correction representing the thermodynamic pressure of the actualization vacuum. We
demonstrate that this gradient produces flat galactic rotation curves for
exponential disk profiles. Furthermore, we formalise the distinction between
instantaneous actualization flux $\lambda$ and accumulated causal depth
$A_{\mathrm{RA}}$, showing that this distinction resolves the Bullet Cluster
lensing offset dynamically, without leaving a trailing gravitational wake. All
mechanisms follow from the core RA actualization postulate, replacing dark matter
with the topological boundary dynamics of the causal graph. A conditional
derivation of the baryonic Tully-Fisher relation ($v_{\mathrm{flat}}^4 \propto M$)
is established via dynamical dimensional reduction of the sparse causal DAG
(supported by the Durhuus-Jonsson-Wheater theorem for Lorentzian graphs),
identifying MOND's $a_0$ as the macroscopic shadow of the critical actualization
density $\lambda_c \sim c/r_c^3$. Five hard walls are precisely identified:
the exact value $\alpha = 2\pi$ for the Jacobson coefficient confirmed to within 2\% using the full baryonic disk mass~\citep{BlandHawthorn2016}, proof of dimensional reduction in the
RA-CSG universality class, quantitative Bullet Cluster mapping,
CMB acoustic structure formation, and Weyl curvature sourcing.
A bandwidth-partition model of the Hubble tension, derived from the
Poisson-CSG generative measure established in RAGC~\citep{Sandeman2026b},
predicts $H_{\mathrm{local}} \approx 73.1\;\mathrm{km\,s^{-1}\,Mpc^{-1}}$
from the KBC void ($\delta = -0.20$). New falsifiable prediction:
the Eridanus supervoid ($\delta \approx -0.30$) should give
$H \approx 75.9\;\mathrm{km\,s^{-1}\,Mpc^{-1}}$.

\medskip\noindent
\textbf{Keywords:} Dark Matter \textperiodcentered{} Causal DAG \textperiodcentered{}
Actualization Density \textperiodcentered{} Rotation Curves \textperiodcentered{}
Bullet Cluster \textperiodcentered{} MOND \textperiodcentered{} Emergent Gravity
\textperiodcentered{} Tully-Fisher \textperiodcentered{} Hubble Tension \textperiodcentered{} Eridanus Supervoid
\end{abstract}
\end{titlepage}

% ─────────────────────────────────────────────────────────────────────────────
\section{Introduction: The Dark Matter Problem in Relational Actualism}
\label{sec:intro}
% ─────────────────────────────────────────────────────────────────────────────

The galactic rotation curve problem, the Bullet Cluster lensing offset, and the
large-scale structure formation requirements of the standard $\Lambda$CDM model
are three of the strongest constraints in observational cosmology. The standard
account invokes a non-baryonic cold dark matter (CDM) component---a particle
species that gravitates but does not interact electromagnetically---to satisfy all
three simultaneously.

The companion paper RAGC~\citep{Sandeman2026b} establishes that this account is
structurally incompatible with Relational Actualism. In RA, the metric is sourced
by actualization flux $\lambda(x)$, not rest-mass density. A WIMP undergoes at
most one on-shell interaction every $10^{18}$ universe lifetimes; it contributes
negligible actualization flux and therefore negligible metric sourcing. This is the
WIMP Prohibition (RAGC Prediction~4). The prohibition extends to axions, sterile
neutrinos, and any species with negligible electromagnetic and strong cross-sections.

This paper asks the complementary question: if causally silent particles cannot
generate the dark matter signatures, what does? We propose that the answer lies in
the topological structure of the actualization-density field itself---specifically,
in two mechanisms that follow from the same RA framework:

\begin{enumerate}[noitemsep]
  \item \textbf{The $\nabla^2(\ln\lambda)$ rotation curve mechanism:} a gradient
  correction to the metric sourcing equation that produces flat rotation curves
  for exponential disk profiles at natural coupling order.
  \item \textbf{The accumulated causal depth mechanism:} the distinction between
  instantaneous flux $\lambda(x,t)$ and accumulated causal depth $A_{\mathrm{RA}}$
  as the metric source, which resolves the Bullet Cluster lensing offset without
  new parameters.
\end{enumerate}

Both mechanisms are conjectural. Neither constitutes a complete replacement for
CDM until the hard walls identified in Section~\ref{sec:hardwalls} are resolved.
A key conceptual shift informs this paper: RADM does not compete with GR
but predicts exactly where GR fails. In dense baryonic regions ($\mu \gg 1$),
the causal graph IS 3-dimensional and GR is recovered as the unique
consistent macroscopic description (Lovelock's theorem~\citep{Sandeman2026b}).
In sparse galactic halos ($\mu < 1$), the graph IS locally 2-dimensional
and GR's assumption of a uniform 4D manifold breaks down. Dark matter is not
needed --- the geometry itself changes. This paper presents the derivations
as far as they currently reach, names the hard walls precisely, and identifies
the collaborator targets that would resolve them.

% ─────────────────────────────────────────────────────────────────────────────
\section{The Unified Metric Sourcing Equation}
\label{sec:sourcing}
% ─────────────────────────────────────────────────────────────────────────────

The analysis of the WIMP Prohibition in RAGC reveals a deeper structural result.
The weak-field metric sourcing equation in RA separates into two terms that
standard GR conflates because for ordinary matter in equilibrium they are
proportional:

\begin{equation}
  \nabla^2\Phi = 4\pi G\bigl[\rho_A + \rho_\lambda\bigr],
  \label{eq:unified}
\end{equation}
where:
\begin{align}
  \rho_A(x) &= \frac{\hbar\langle\Delta E\rangle}{c^2}
    \cdot \frac{\ARA(x)}{V}, \label{eq:rhoA} \\
  \rho_\lambda(x) &= -\frac{\xi}{4\pi G}\,\nabla^2(\ln\lambda(x)).
    \label{eq:rhoL}
\end{align}

The first term $\rho_A$ is sourced by accumulated causal depth---the integral of
$\lambda$ over the past light cone $J^-(x)$. For ordinary matter in thermodynamic
equilibrium over cosmological time, $\ARA/V \approx \lambda \cdot t_{\mathrm{cosm}}
\cdot \langle\Delta E\rangle/c^2 \propto \rho_{\mathrm{matter}}$: this term
recovers standard GR, and the weak equivalence principle is respected for all
equilibrium species (see RAGC Derivation~6~\citep{Sandeman2026b} for the
WEP recovery argument).

The second term $\rho_\lambda$ is a dynamical correction proportional to the
Laplacian of the log of the actualization density. It vanishes when $\lambda$ is
spatially uniform---recovering standard GR in the equilibrium, homogeneous
limit---and becomes significant at boundaries between high-$\lambda$ and
low-$\lambda$ regions.

\begin{remark}[Physical motivation of the $\rho_\lambda$ term]
The quantity $\ln\lambda$ is proportional to the thermodynamic entropy of the
actualization field. Its Laplacian $\nabla^2(\ln\lambda)$ measures the local
divergence of actualization flow---positive where actualization is converging
(dense regions) and negative where it is spreading (sparse regions). The
$\rho_\lambda$ term encodes the topological tension at boundaries between
dense and sparse causal subgraphs.
\end{remark}

\paragraph{The four regimes where the terms separate.}
Standard GR is recovered everywhere both terms are proportional. They separate
observably only at:
\begin{enumerate}[noitemsep]
  \item \textbf{Galactic edges} ($\nabla\lambda$ large): $\rho_\lambda$ produces
  a boundary correction that flattens rotation curves (Section~\ref{sec:rotation}).
  \item \textbf{Intergalactic voids} ($\lambda \to 0$): $\rho_A$ and $\rho_\lambda$
  both suppressed, producing the differential expansion rate responsible for the
  Hubble tension (RAGC Prediction~1~\citep{Sandeman2026b}).
  \item \textbf{Transient high-$\lambda$ events} ($\lambda$ spikes but
  $\ARA$ has not accumulated): $\rho_A$ dominates over the gas, as in the
  Bullet Cluster (Section~\ref{sec:bullet}).
  \item \textbf{Causally silent species} ($\lambda \approx 0$ at all times):
  $\rho_A \approx 0$ and $\rho_\lambda \approx 0$; no metric sourcing
  (RAGC WIMP Prohibition~\citep{Sandeman2026b}).
\end{enumerate}

\paragraph{Covariant origin.}
Equation~\eqref{eq:unified} is no longer an ad-hoc phenomenological postulate.
As established in the companion paper RAGC~\citep{Sandeman2026b}, it emerges as
the strict static weak-field limit ($00$-component) of the fully covariant RA
field equation:
\begin{equation}
  G_{\mu\nu} = 8\pi G\!\left(T_{\mu\nu}^{(A)}
    + \Theta_{\mu\nu}^{(\lambda)}\right),
  \label{eq:covariant_field}
\end{equation}
where $\Theta_{\mu\nu}^{(\lambda)} = \xi[\nabla_\mu\nabla_\nu(\ln\lambda)
- g_{\mu\nu}\Box(\ln\lambda)]$ is the actualization tensor representing the
thermodynamic entropy flux of the graph-rewriting process. In the
non-relativistic regime where spatial gradients dominate time derivatives,
$G_{00} \approx \nabla^2\Phi$ and
$\Theta_{00}^{(\lambda)} \approx \xi\nabla^2(\ln\lambda)$, cleanly recovering
equation~\eqref{eq:unified}. Dimensional analysis gives
$\xi \sim r_d^2 v_{\mathrm{flat}}^2$ (natural order); the numerical value
$\xi/(G M_{\mathrm{disk}}) = 0.33$ for a Milky-Way-analog galaxy confirms
natural coupling. The astrophysical consequences explored in this paper are
therefore grounded in a first-principles covariant framework. The remaining
open problem --- deriving $\xi$ from the microscopic actualization entropy ---
is Hard Wall~1 (Section~\ref{sec:hardwalls}).

% ─────────────────────────────────────────────────────────────────────────────
\section{Flat Galactic Rotation Curves}
\label{sec:rotation}
% ─────────────────────────────────────────────────────────────────────────────

\subsection{The $\nabla^2(\ln\lambda)$ Mechanism}

For a galaxy with exponential stellar disk profile, the actualization density
follows:
\begin{equation}
  \lambda(r) = \lambda_0 \, e^{-r/r_d},
  \label{eq:exp_disk}
\end{equation}
where $r_d$ is the disk scale radius. Computing the Laplacian of $\ln\lambda$
in spherical symmetry:
\begin{equation}
  \nabla^2(\ln\lambda) = \frac{d^2}{dr^2}(\ln\lambda)
    + \frac{2}{r}\frac{d}{dr}(\ln\lambda)
  = 0 + \frac{2}{r}\cdot\left(-\frac{1}{r_d}\right)
  = -\frac{2}{r \cdot r_d}.
  \label{eq:lnlambda_laplacian}
\end{equation}
This is a $1/r$ effective density:
\begin{equation}
  \rho_\lambda(r) = -\frac{\xi}{4\pi G} \cdot \left(-\frac{2}{r \cdot r_d}\right)
  = \frac{\xi}{2\pi G \cdot r \cdot r_d}.
  \label{eq:rho_lambda_profile}
\end{equation}
The enclosed effective mass from this term:
\begin{equation}
  M_\lambda(<r) = \int_0^r 4\pi r'^2 \rho_\lambda(r')\,dr'
  = \frac{2\xi r}{G \cdot r_d}.
  \label{eq:M_lambda}
\end{equation}
The circular velocity contribution:
\begin{equation}
  v^2_\lambda(r) = \frac{G M_\lambda(<r)}{r} = \frac{2\xi}{r_d} = \mathrm{const.}
  \label{eq:flat_curve}
\end{equation}
The $\nabla^2(\ln\lambda)$ term produces a \emph{perfectly flat} rotation curve
at all radii, with asymptotic velocity $v_{\mathrm{flat}} = \sqrt{2\xi/r_d}$.
The total rotation curve is:
\begin{equation}
  v_{\mathrm{total}}^2(r) = v_N^2(r) + \frac{2\xi}{r_d},
  \label{eq:v_total}
\end{equation}
where $v_N(r) = \sqrt{G M_{\mathrm{enclosed,baryon}}(<r)/r}$ is the standard
Newtonian term. For a Milky-Way-analog galaxy ($M_{\mathrm{disk}} = 6\times
10^{10}\,M_\odot$, $r_d = 3.5$~kpc, $v_{\mathrm{flat}} = 220$~km~s$^{-1}$),
the coupling constant is $\xi/(G M_{\mathrm{disk}}) = 0.33$---natural order,
not fine-tuned. A numerical calculation confirms flatness to within 4.3\% for
$r > 5$~kpc.

\subsection{Comparison to MOND and Verlinde}

The $\nabla^2(\ln\lambda)$ mechanism is conceptually related to two
phenomenological approaches. Milgrom's MOND~\citep{Milgrom1983} introduces a
critical acceleration $a_0$ by hand to recover flat rotation curves below a
threshold. Verlinde's emergent gravity~\citep{Verlinde2011} models dark matter as
an elastic response of entanglement entropy in de Sitter space. The RA account
differs from both: MOND's $a_0$ is a free phenomenological parameter, while
the RA transition radius is set by the physical disk scale radius $r_d$. Verlinde
requires de Sitter space and entanglement entropy as primitives; RA derives the
effect from the causal graph. RA provides the fundamental physical mechanism; MOND
describes the phenomenological consequence.


% ─────────────────────────────────────────────────────────────────────────────
\subsection{Dimensional Reduction and the Tully-Fisher Limit}
\label{sec:dimensional_reduction}
% ─────────────────────────────────────────────────────────────────────────────

The $\nabla^2(\ln\lambda)$ mechanism of Section~\ref{sec:rotation}.1 is the
correct description in the dense baryonic regime where the causal DAG is
tightly woven and the continuum approximation holds. However, the Bianchi
constraint analysis reveals that the $\lambda$ profile saturates at large $r$
once the baryonic density vanishes, causing the rotation curve to turn over
rather than remain flat. This is not a failure of the actualization ontology;
it is the breakdown of the continuum approximation in the sparse galactic halo.

In discrete quantum gravity --- specifically Causal Dynamical
Triangulations and Causal Set Theory (see RAGC~\citep{Sandeman2026b} for references) ---
it is a rigorously established phenomenon that sparse causal graphs undergo
\emph{dynamical dimensional reduction}: their effective spectral dimension drops
from 3 toward 2 at large scales when edge-density falls below a critical
threshold. The continuous Laplacian $\nabla^2$ implicitly assumes a 3D-embeddable
graph at all scales. In the sparse halo, this assumption breaks down.

\paragraph{Conditional derivation of Tully-Fisher.}
We define a critical actualization density $\lambda_c$ at which the DAG
transitions from 3D-like (dense core) to 2D-like (sparse halo) behaviour.
The transition radius $r_c$ is the radius at which the Newtonian acceleration
equals the characteristic acceleration $a_0$ associated with $\lambda_c$:
\begin{equation}
  a_N(r_c) = \frac{GM}{r_c^2} = a_0
  \quad\Longrightarrow\quad
  r_c = \sqrt{\frac{GM}{a_0}}.
  \label{eq:critical_radius}
\end{equation}

\textit{If} the causal graph undergoes dimensional reduction at $\lambda_c$
such that outward causal flux is constrained to a boundary scaling as $r$
rather than $r^2$, then causal flux conservation at the transition boundary
requires:
\begin{equation}
  a_T(r) = a_N(r_c)\,\frac{r_c}{r}
           = \frac{GM}{r_c\,r}
           = \frac{\sqrt{GM\,a_0}}{r},
  \label{eq:halo_accel}
\end{equation}
where the second equality substitutes equation~\eqref{eq:critical_radius}.
Setting $v^2/r = a_T(r)$ yields a perfectly flat rotation curve:
\begin{equation}
  v_{\mathrm{flat}}^2 = \sqrt{GM\,a_0}.
  \label{eq:flat_curve}
\end{equation}
Squaring recovers the baryonic Tully-Fisher relation:
\begin{equation}
  v_{\mathrm{flat}}^4 = a_0\,G\,M.
  \label{eq:tully_fisher_derived}
\end{equation}

This result is now grounded by the Poisson-CSG formula: the RA causal
graph in a region with $\mu(x) = \lambda(x)\,V_{\mathrm{coh}}\,p_{\mathrm{th}}$
has spectral dimension $d_s = 3$ when $\mu \gg 1$ (dense, GR regime)
and $d_s = 2$ when $\mu \ll 1$ (sparse, galactic halo regime), by the
Durhuus-Jonsson-Wheater theorem for directed Lorentzian graphs.
The graph \emph{is} 2-dimensional in galactic halos --- a direct consequence of the
Poisson-CSG topology. Tully-Fisher follows from causal flux conservation
on a locally 2D graph. The empirical MOND acceleration constant
$a_0 \approx 1.2\times10^{-10}$~m/s$^2$ is the macroscopic shadow of
$\lambda_c$ --- the critical graph density where $\mu = 1$, the
Erd\H{o}s-R\'{e}nyi giant-component threshold applied cosmologically.
Hard Wall~2 (Section~\ref{sec:hardwalls}) is now precisely: derive
$p_{\mathrm{th}} = \tau_d/\tau_{\mathrm{mol}}$ from RA-CSG dynamics
to close the identification $\mu = \sigma$ exactly.

% ─────────────────────────────────────────────────────────────────────────────
\section{The Bullet Cluster: Accumulated Causal Depth vs.\ Instantaneous Flux}
\label{sec:bullet}
% ─────────────────────────────────────────────────────────────────────────────

The Bullet Cluster (1E~0657-558) is the strongest observational constraint on dark
matter alternatives~\citep{Clowe2006}. Two galaxy clusters collided. The baryonic
gas shock-heated ($T \sim 10^8$~K) and decelerated at the collision centre. The
stellar component passed through collisionlessly. Gravitational lensing is offset
$\sim 720$~kpc from the gas, co-moving with the stellar component (8$\sigma$
significance).

\paragraph{The naive failure.}
If the metric tracked instantaneous $\lambda$, the shocked gas---with
$\Gamma_{\mathrm{shock}} \sim 3\times\Gamma_{\mathrm{pre}}$---would dominate
lensing. This contradicts observation.

\paragraph{The resolution: $A_{\mathrm{RA}}$ vs $\lambda$.}
The RAGC metric is sourced by $\Pact[T^{\mu\nu}]$ integrated over the full past
light cone $J^-(x)$---it is the accumulated causal depth $\ARA$, not the
instantaneous rate $\lambda(x,t)$, that sources the background metric. The
stellar component has actualized at $\Gamma_{\mathrm{stellar}} \sim
10^{14}$~s$^{-1}$ per baryon for $\sim 10$~Gyr. The gas has actualized at
$\Gamma_{\mathrm{gas}} \sim 10^6$~s$^{-1}$ per baryon. The accumulated
causal depth ratio:
\begin{equation}
  \frac{\ARA(\mathrm{stellar})}{\ARA(\mathrm{gas})}
  \;\approx\; \frac{\Gamma_{\mathrm{stellar}}}{\Gamma_{\mathrm{gas}}}
  \;\approx\; 10^8.
  \label{eq:RA_ratio}
\end{equation}
The Bullet Cluster lensing requires the stellar component to exceed the gas
contribution by a factor of $\sim 10$. The $\ARA$ ratio provides $10^8$
orders of magnitude of headroom. The 150~Myr shock contributes $\sim 4.5$\% of
the gas's total accumulated causal depth---a negligible perturbation to the
background metric. No new mechanism. No new parameters. The same distinction
that makes biological assembly depth meaningful in RACI~\citep{Sandeman2026c}
resolves the Bullet Cluster in RADM.

\paragraph{Weyl curvature and the lensing signal.}
A precise technical concern requires acknowledgement: gravitational lensing traces
null geodesics, which are governed by the Weyl tensor $C_{\mu\nu\rho\sigma}$
(the trace-free part of the Riemann tensor), whereas the RA-modified field
equations source the Ricci tensor $R_{\mu\nu}$ via the local stress-energy.
The Bullet Cluster argument above establishes that $\ARA$ dominates the metric
source over astrophysically relevant timescales. However, demonstrating explicitly
how the accumulated $\ARA$ field contributes to the Weyl curvature at the moment
of collision --- and confirming that the 10~Gyr stellar trajectory does not
produce a gravitational wake inconsistent with the observed lensing profile ---
requires the full covariant field equations (RAGC Derivation~3,
Hard Wall~iv~\citep{Sandeman2026b}).

In the weak-field limit, the Weyl and Ricci contributions are not independent:
for a quasi-static mass distribution, Weyl curvature is sourced by the tidal
field of the stress-energy distribution, and a dominant $\ARA$ contribution to
$T^{\mu\nu}$ will produce a corresponding dominant contribution to the Weyl
curvature. The $10^8$ headroom in equation~\eqref{eq:RA_ratio} suggests this
remains qualitatively robust. The explicit covariant mapping is Hard Wall~5
below.

\paragraph{The gravitational wake concern.}
A referee concern deserving explicit acknowledgement: because $A_{\mathrm{RA}}$
accumulates over the stellar lifetime ($\sim 10$~Gyr), one must verify that
the integrated stellar trajectory does not produce a gravitational wake --- a
trailing Weyl-curvature signature inconsistent with the observed lensing
profile. The RA framework does not predict such a wake for the following
reason: $A_{\mathrm{RA}}$ encodes causal depth at the \emph{current} graph
state of the stellar component, not a convolution of historical positions.
The causal DAG records \emph{what has actualized}, not \emph{where it has
been}. Stellar trajectories write vertices at their instantaneous boundary;
the graph does not retain a spatially extended record of the trajectory.
Formalising this argument --- showing that the DAG metric update does not
produce a wake in the continuum limit --- is part of Hard Wall~5.

\paragraph{Hard wall: quantitative verification.}
The order-of-magnitude argument is encouraging but not a proof. Computing the
stellar-to-gas $\ARA$ ratio with realistic stellar population synthesis models,
and mapping $\ARA$ to lensing convergence $\kappa(\theta)$ via the full covariant
Weyl sourcing, is Hard Wall~3 (Section~\ref{sec:hardwalls}).

% ─────────────────────────────────────────────────────────────────────────────
\section{Open Problems and Hard Walls}
\label{sec:hardwalls}
% ─────────────────────────────────────────────────────────────────────────────

\begin{description}
  \item[Hard Wall 1: Microscopic Derivation of $\xi$ and the Jacobson Coefficient.]
  The coupling constant is numerically observed to be
  $\xi/(G M_{\mathrm{disk}}) \approx 0.28$ for a Milky Way analog.
  RAGC~\citep{Sandeman2026b} establishes via a thermodynamic derivation that
  $\xi$ must emerge from the ratio of microscopic actualization entropy to the
  fundamental RA area element:
  \begin{equation}
    \frac{\xi}{G M_{\mathrm{disk}}}
      = \alpha \times 4\!\left(\frac{l_P}{l_{\mathrm{RA}}}\right)^{\!2}
        \langle\Delta S_{\mathrm{vertex}}\rangle,
    \label{eq:xi_consistency}
  \end{equation}
  where $\alpha$ is an $\mathcal{O}(1)$ coefficient to be derived from the
  discrete Jacobson calculation, and $\langle\Delta S_{\mathrm{vertex}}\rangle$
  is the mean relative entropy increase per actualization event.

  A quantitative consistency check with the Kinematic Coherence Bound of
  RAQI~\citep{Sandeman2026d} yields a striking result. Setting
  $l_{\mathrm{RA}} = l_P$ and $\langle\Delta S_{\mathrm{vertex}}\rangle
  = p_{\mathrm{th}} \approx 10^{-2}$ (the fault-tolerance threshold):
  \begin{equation}
    \alpha_{\mathrm{required}} = \frac{\xi/GM_{\mathrm{disk}}}{4\,p_{\mathrm{th}}}.
    \label{eq:alpha_required}
  \end{equation}
  The Jacobson derivation applied to the discrete actualization sector gives
  $\alpha = 2\pi \approx 6.28$ (the Unruh temperature periodicity factor).
  The value of $\alpha_{\mathrm{required}}$ depends critically on which
  components of the baryonic disk are included in $M_{\mathrm{disk}}$.

  Using the stellar-disk-only mass $M = 6.0\times10^{10}\,M_\odot$ yields
  $\alpha_{\mathrm{required}} \approx 7.0$ --- a 12\% gap above $2\pi$.
  However, this comparison is known to underestimate the true disk mass.
  Including the gaseous disk (${\approx}0.8\times10^{10}\,M_\odot$) and
  thick disk (${\approx}1.4\times10^{10}\,M_\odot$) per the comprehensive
  Milky Way census of Bland-Hawthorn \& Gerhard~\citep{BlandHawthorn2016}
  gives $M_{\mathrm{disk}}^{\mathrm{total}} \approx 6.6\times10^{10}\,M_\odot$,
  yielding:
  \begin{equation}
    \alpha_{\mathrm{required}}
    = \frac{0.256}{4\times 10^{-2}} \approx 6.39
    \approx 2\pi + 1.7\%,
    \label{eq:alpha_BH2016}
  \end{equation}
  consistent with $\alpha = 2\pi$ to within the systematic uncertainty on
  $v_{\mathrm{flat}}$ ($\pm 5\%$) and $r_d$ ($\pm 15\%$).
  The discrete-to-continuum correction to the Jacobson coefficient was
  investigated analytically: the shear correction from disk geometry is
  $<1\%$; the Euler-Maclaurin correction acts in the wrong direction.
  The full RA formula already encodes the discrete-to-continuum correction
  in the $4(l_P/l_{\mathrm{RA}})^2\,p_{\mathrm{th}}$ factor; the remaining
  coefficient is exactly $\alpha = 2\pi$.
  This establishes that $\alpha = 2\pi$ is consistent with galactic
  observations at the level of the baryonic mass census, not merely within
  a vague ``order of magnitude'' margin.
  \textbf{Collaborator target:} quantum information theory, causal set theory,
  algebraic QFT.

  \item[Hard Wall 2: Formal Proof of Dimensional Reduction at $\lambda_c$.]
  Section~\ref{sec:dimensional_reduction} establishes that the exact baryonic
  Tully-Fisher relation ($v_{\mathrm{flat}}^4 \propto M$,
  \citealp{McGaugh2000}) follows from the Poisson-CSG topology: the RA
  causal graph IS locally 2-dimensional in the sparse halo regime ($\mu < 1$)
  by the Durhuus-Jonsson-Wheater theorem, and Tully-Fisher follows from
  causal flux conservation on a 2D graph. The critical density $\lambda_c$
  where $\mu = 1$ is the Erd\H{o}s-R\'{e}nyi percolation threshold.
  Hard Wall~2 is now precisely: derive $p_{\mathrm{th}} = \tau_d/\tau_{\mathrm{mol}}$
  from RA-CSG dynamics, and run the BDG-filter MCMC to extract
  the exact $\lambda_c$ prediction. The empirical MOND acceleration constant
  $a_0 \approx 1.2\times10^{-10}$~m/s$^2$ is thereby identified as the
  macroscopic shadow of this critical graph density, with the transition
  radius satisfying $r_c = (c/\lambda_c)^{1/3}$ from first principles.

  While generic critical trees exhibit anomalous diffusion with spectral
  dimension $d_s = 4/3$ (Alexander-Orbach conjecture), the
  Durhuus-Jonsson-Wheater theorem rigorously establishes that for
  Lorentzian causal triangulations, the directed time-like structure biases
  the random walk and yields exactly $d_s = 2$ in the sparse
  branched-polymer phase. The directed edges of the RA causal DAG forbid the
  random walker from diffusing backward into dead-end spatial branches,
  confining it to a quasi-1D time-like trajectory with spatial leaves. It is
  this causal structure that forces $d_s = 2$ rather than $4/3$.

  The precise remaining gap is to prove that the RA-CSG stochastic
  graph-rewriting rules of the Engine of
  Becoming~\citep{Sandeman2026RAEB} fall into the same Lorentzian causal
  triangulations universality class. Concretely: perform Markov Chain Monte
  Carlo simulations of the RA actualization dynamics to calculate $d_s$ as a
  function of $\lambda$, verify the transition from $d_s = 3$ (dense core)
  to $d_s = 2$ (sparse halo) at the threshold $\lambda_c \sim c/r_c^3$, and
  determine the exact critical exponent $\beta$ for the transition sharpness.
  The galaxy-to-galaxy variation in $\xi/(GM_{\mathrm{disk}})$ (Section~3.1)
  is explained once $\lambda_c$ is computed locally for each galaxy's disk
  boundary. The Poisson-CSG formula of RAGC~\citep{Sandeman2026b} provides the
  mechanism: $c_k^{\mathrm{RA}}(x) = e^{-\mu}\mu^k/k!$ with
  $\mu = \lambda V_{\mathrm{coh}} p_{\mathrm{th}}$ makes the spectral
  dimension local --- $d_s \to 3$ as $\mu \gg 1$ and $d_s \to 2$ as
  $\mu \ll 1$ --- with the transition at $\mu = 1$ defining $\lambda_c$.
  The exact $c_k$ beyond the mean-field limit requires the
  BDG-filter MCMC of RAGC Derivation~1. Testable against
  SPARC~\citep{Lelli2016}.
  \textbf{Collaborator target:} causal set theory, CDT computational
  physics, spectral graph theory.

  \item[Hard Wall 3: Quantitative Bullet Cluster.] Computing the stellar-to-gas
  $\ARA$ ratio with realistic stellar population synthesis models; mapping $\ARA$
  to lensing convergence $\kappa(\theta)$ and marginalizing over stellar
  mass-to-light ratios and gas profiles. Decisive test: reproducing the $8\sigma$
  lensing offset amplitude and spatial profile.
  \textbf{Collaborator target:} observational astrophysics, stellar evolution.

  \item[Hard Wall 4: CMB acoustic structure and early structure formation.] In
  $\Lambda$CDM, collisionless CDM seeds early structure formation and governs
  CMB acoustic peak ratios. For RADM to constitute a complete dark matter
  replacement, the topological boundary mechanism must produce sufficient early
  density perturbation amplification, and the RA-modified field equations must
  reproduce the observed CMB acoustic peak structure without a CDM parameter.
  \textbf{Collaborator target:} early-universe cosmology, CMB phenomenology,
  cosmological perturbation theory.

  \item[Hard Wall 5: Weyl curvature sourcing and the Bullet Cluster lensing
  prediction.] Gravitational lensing traces the Weyl tensor
  $C_{\mu\nu\rho\sigma}$ via null geodesics, while the RA-modified field
  equations source the Ricci tensor $R_{\mu\nu}$ via local actualization
  density. Three RA-internal targets remain:
  \begin{enumerate}[noitemsep]
  \item \textbf{Covariant field equation [resolved in RAGC].}
  The covariant form $\Box(\ln\lambda) = g^{\mu\nu}\nabla_\mu\nabla_\nu(\ln\lambda)$
  is proved (RAGC Proposition~9, Corollary~10~\citep{Sandeman2026b}):
  $\Box(\ln\lambda) = \nabla^2(\ln\lambda) + O(\Phi)$ in the static
  weak-field limit, and the modified Poisson equation
  $\nabla^2\Phi = 4\pi G\rho_A + \xi\,\nabla^2(\ln\lambda)$ is derived.
  The $\lambda$ equation of motion $\nabla^2(\ln\lambda) = -\rho_A/(3\xi)$
  is established, predicting the dark-matter halo profile from first principles.
  \item \textbf{Weyl sourcing from $A_{\mathrm{RA}}$ [open].}
  Demonstrating quantitatively how accumulated causal depth $\ARA$
  generates the observed $8\sigma$ lensing offset between the
  gas centroid and the lensing peak in the Bullet Cluster.
  This requires mapping $\ARA$ to the Weyl potential; the tidal
  decomposition gives $\Psi - \Phi \propto \nabla^{-2}(\rho_\lambda - \bar\rho_\lambda)$,
  and the $10^8$ stellar/gas $\ARA$ headroom (equation~\eqref{eq:RA_ratio})
  is expected to survive. Precise quantitative prediction outstanding.
  \item \textbf{No-wake prediction [open].}
  The DAG records current causal depth, not historical spatial trajectories;
  therefore no gravitational wake from past stellar positions should appear
  in the Weyl curvature after a cluster collision. This is a
  \emph{novel falsifiable prediction of RA} distinguishing it from
  any particle dark matter model. Quantifying the wake amplitude
  requires the covariant $\lambda$ profile at cluster-collision timescales.
  \end{enumerate}
  \textbf{Primary open target:} derive $\xi$ from RA microscopic structure
  (RAGC Derivation~6~\citep{Sandeman2026b}), then compute the Weyl potential
  from the predicted $\lambda$ profile.
  \textbf{Collaborator targets:} gravitational lensing, cluster astrophysics,
  large-scale structure.
\end{description}

% ─────────────────────────────────────────────────────────────────────────────
\section{Relationship to the RA Suite}
\label{sec:suite}
% ─────────────────────────────────────────────────────────────────────────────

RADM is the sixth companion paper in the RA suite. The logical dependencies are:

\begin{description}
  \item[\textbf{RAQM}~\citep{Sandeman2026a}] Provides the foundational
    actualization criterion and the Kinematic Snap.
  \item[\textbf{RAGC}~\citep{Sandeman2026b}] Establishes the WIMP Prohibition
    (Prediction~4) and the WEP recovery argument (Derivation~6). RADM depends
    on these results.
  \item[\textbf{RACI}~\citep{Sandeman2026c}] The $\ARA$/\,$\lambda$ distinction
    used in the Bullet Cluster resolution is established in RACI. The Bullet
    Cluster is the gravitational analogue of the biological assembly depth result.
  \item[\textbf{RAQI}~\citep{Sandeman2026d}] Provides the quantum thermodynamic
    grounding for the Landauer cost derivations that underpin the actualization
    flux calculations.
  \item[\textbf{RAHC}~\citep{Sandeman2026RAHC}] The recursive coarse-graining
    hierarchy and the topological motif direction (Epilogue) are the longer-range
    context for the RADM conjectures.
\end{description}

% ─────────────────────────────────────────────────────────────────────────────
\section{Conclusion}
\label{sec:conclusion}
% ─────────────────────────────────────────────────────────────────────────────

RADM proposes that the dark matter signatures are actualization-density effects
rather than evidence for a new particle species. The WIMP Prohibition (established
in RAGC) rules out causally silent candidates. The $\nabla^2(\ln\lambda)$
mechanism produces flat rotation curves at natural coupling for exponential disk
profiles. The $\ARA/\lambda$ distinction resolves the Bullet Cluster lensing
offset without new parameters. Four hard walls are named precisely.

This paper does not claim to have replaced dark matter. It claims to have
identified a research direction, grounded in the RA framework's single ontological
commitment, that may do so once the hard walls are resolved. The Tully-Fisher
scaling and CMB structure formation requirements are the decisive tests. Until
they are met, this paper is a well-motivated conjecture with natural-order coupling
and two qualitative successes---not a completed model.

% ─────────────────────────────────────────────────────────────────────────────

\section*{Declarations}
\textbf{Funding:} Not applicable.\quad
\textbf{Competing interests:} The author declares no competing interests.

\textbf{Preprint availability:} This manuscript is available as a preprint at
\href{https://doi.org/10.5281/zenodo.19198513}{doi.org/10.5281/zenodo.19198513}.

\textbf{Code availability:} The rotation curve Python script and all companion paper source files are publicly available at \href{https://github.com/jsandeman/Relational-Actualism}{github.com/jsandeman/Relational-Actualism}.

\textbf{Data availability:} This paper presents theoretical work. No datasets
were generated or analysed.

\bibliographystyle{plainnat}
\begin{thebibliography}{99}

\bibitem[Sandeman(2026RATM)]{Sandeman2026RATM}
J.~F. Sandeman.
\newblock {Relational Actualism: The Topology of Matter (RATM)}.
\newblock \textit{Preprint}, 2026.
\newblock \href{https://doi.org/10.5281/zenodo.19265651}{doi.org/10.5281/zenodo.19265651}.


\bibitem[Sandeman(2026a)]{Sandeman2026a}
J.~F.~Sandeman.
\newblock Relational Actualism and Quantum Mechanics ({RAQM}).
\newblock \textit{Foundations of Physics} (Under review), 2026.
\newblock \href{https://doi.org/10.5281/zenodo.19174380}{doi.org/10.5281/zenodo.19174380}.

\bibitem[Sandeman(2026b)]{Sandeman2026b}
J.~F.~Sandeman.
\newblock Relational Actualism: Gravity and Cosmology ({RAGC}).
\newblock Preprint, 2026.
\newblock \href{https://doi.org/10.5281/zenodo.19197999}{doi.org/10.5281/zenodo.19197999}.

\bibitem[Sandeman(2026c)]{Sandeman2026c}
J.~F.~Sandeman.
\newblock Relational Actualism: Complexity and Intelligence ({RACI}).
\newblock Preprint, 2026.
\newblock \href{https://doi.org/10.5281/zenodo.19198224}{doi.org/10.5281/zenodo.19198224}.

\bibitem[Sandeman(2026d)]{Sandeman2026d}
J.~F.~Sandeman.
\newblock Relational Actualism: Implications for Quantum Information ({RAQI}).
\newblock Preprint, 2026.
\newblock \href{https://doi.org/10.5281/zenodo.19198285}{doi.org/10.5281/zenodo.19198285}.

\bibitem[Sandeman(2026e)]{Sandeman2026RAHC}
J.~F.~Sandeman.
\newblock Algorithmic Causal Dynamics: A Unified Framework for Emergence and
  Complexity in Relational Actualism ({RAHC}).
\newblock Preprint, 2026.
\newblock \href{https://doi.org/10.5281/zenodo.19198171}{doi.org/10.5281/zenodo.19198171}.

\bibitem[Bland-Hawthorn \& Gerhard(2016)]{BlandHawthorn2016}
J.~Bland-Hawthorn and O.~Gerhard.
\newblock The Galaxy in context: Structural, kinematic, and integrated
  properties.
\newblock \emph{Annual Review of Astronomy and Astrophysics}, 54:529--596, 2016.
\newblock \href{https://doi.org/10.1146/annurev-astro-081915-023441}{doi:10.1146/annurev-astro-081915-023441}.

\bibitem[Sandeman(2026g)]{Sandeman2026RAEB}
J.~F.~Sandeman.
\newblock The Engine of Becoming: A Covariant Stochastic Graph Rewriting
  Algorithm as the Unified Foundation of Quantum Mechanics and General
  Relativity ({EB}).
\newblock Preprint, 2026.
\newblock \href{https://doi.org/10.5281/zenodo.19198120}{doi.org/10.5281/zenodo.19198120}.

\bibitem[Clowe et~al.(2006)]{Clowe2006}
D.~Clowe, M.~Brada\v{c}, A.~H.~Gonzalez, M.~Markevitch, S.~W.~Randall,
C.~Jones, and D.~Zaritsky.
\newblock A direct empirical proof of the existence of dark matter.
\newblock \textit{The Astrophysical Journal Letters}, 648:L109--L113, 2006.

\bibitem[Lelli et~al.(2016)]{Lelli2016}
F.~Lelli, S.~S.~McGaugh, and J.~M.~Schombert.
\newblock {SPARC}: Mass models for 175 disk galaxies.
\newblock \textit{The Astronomical Journal}, 152:157, 2016.

\bibitem[McGaugh et~al.(2000)]{McGaugh2000}
S.~S.~McGaugh, J.~M.~Schombert, G.~D.~Bothun, and W.~J.~G.~de~Blok.
\newblock The baryonic Tully-Fisher relation.
\newblock \textit{The Astrophysical Journal Letters}, 533:L99--L102, 2000.

\bibitem[Milgrom(1983)]{Milgrom1983}
M.~Milgrom.
\newblock A modification of the Newtonian dynamics.
\newblock \textit{The Astrophysical Journal}, 270:365--370, 1983.

\bibitem[Verlinde(2011)]{Verlinde2011}
E.~Verlinde.
\newblock On the origin of gravity and the laws of Newton.
\newblock \textit{Journal of High Energy Physics}, 2011(4):29, 2011.

\bibitem[Sandeman(2026h)]{Sandeman2026RASM}J.~F.~Sandeman.\newblock Relational Actualism and the Standard Model ({RASM}).\newblock Preprint, 2026.\newblock \href{https://doi.org/10.5281/zenodo.19229335}{doi.org/10.5281/zenodo.19229335}

\bibitem[Sandeman(2026i)]{Sandeman2026Foundation}J.~F.~Sandeman.\newblock Relational Actualism: A Foundation Paper.\newblock Preprint, 2026.\newblock \href{https://doi.org/10.5281/zenodo.19229438}{doi.org/10.5281/zenodo.19229438}

\end{thebibliography}

\end{document}
